\documentclass[11pt]{article}
\usepackage[margin=1in]{geometry}
\usepackage{amsmath,amssymb,bm,booktabs,array,microtype}
\usepackage[colorlinks=true,allcolors=blue]{hyperref}
\usepackage[nameinlink,noabbrev]{cleveref}
\newcommand{\KK}{\mathcal{K}}
\newcommand{\Rcal}{\mathcal{R}}
\newcommand{\dd}{\mathrm{d}}
\newcommand{\eff}{\mathrm{eff}}
\title{Diffusion-Limited Reaction Frequencies for One-Dimensionally Migrating Defect Clusters}
\author{}
\date{}
\begin{document}
\maketitle
\begin{abstract}
This note derives reaction frequencies for a one-dimensionally (1D) migrating
defect cluster interacting with fixed sinks, extends the result to 1D motion
with stochastic changes of glide direction (1DR--0), and develops the three
geometrical cases for reactions between two purely 1D-mobile clusters. The
notation is chosen for cluster-dynamics equations written in atomic fractions.
Particular care is taken to distinguish a volumetric bimolecular coefficient,
a sink strength, a loss frequency in hertz, and a total reaction-rate density.
\end{abstract}

\section{Definitions and normalization}
For BCC iron,
\begin{equation}
 \Omega=\frac{a^3}{2},\qquad c_\alpha=\Omega N_\alpha,\qquad
 D_\alpha=a^2\omega_\alpha .
\end{equation}
Here \(N_\alpha\) is a number density, \(c_\alpha\) is an atomic fraction, and
\(\omega_\alpha\) has units s\(^{-1}\). Four related quantities must not be
confused:
\begin{align}
 \dot N_A&=-K^{\rm vol}_{AB}N_AN_B,
 &[K^{\rm vol}]&={\rm m^3\,s^{-1}},\\
 \dot c_A&=-K^{\rm af}_{AB}c_Ac_B,
 &K^{\rm af}_{AB}&=K^{\rm vol}_{AB}/\Omega\quad [{\rm s^{-1}}],\\
 \dot N_A&=-D_A k_A^2N_A,
 &[k_A^2]&={\rm m^{-2}},\\
 \dot c_A&=-\KK_Ac_A,
 &\KK_A&=D_Ak_A^2\quad [{\rm s^{-1}}].
 \label{eq:definitions}
\end{align}
In this note, \emph{kernel in hertz} means the loss frequency \(\KK_A\).
Such a frequency may depend on the instantaneous sink and reactant
populations; it is not necessarily a concentration-independent binary
constant.

\section{Pure 1D migration to fixed sinks}
\subsection{Transport analogy: microscopic and macroscopic cross sections}
Consider a cluster of class \(n\) gliding along a fixed line. A fixed sink of
family \(m\) captures it when the glide line passes through the sink's
projected capture area. In analogy with neutron transport, define
\begin{equation}
 \sigma_{n,m}=\pi R_{n,m}^2                         \label{eq:sigma}
\end{equation}
as the microscopic capture cross section. The macroscopic cross sections are
\begin{equation}
 \Sigma_{n,m}=\sigma_{n,m}N_m,\qquad
 \Sigma_n=\sum_j\Sigma_{n,j}=\sum_j\sigma_{n,j}N_j,
 \label{eq:Sigma}
\end{equation}
and the geometrical mean free path along an uncorrelated line is
\begin{equation}
 \lambda_{g,n}=\Sigma_n^{-1}.                       \label{eq:mfp}
\end{equation}
Thus \(\sigma\) has units m\(^{2}\), \(\Sigma\) has units m\(^{-1}\), and
\(\lambda_g\) has units m. The symbol \(\kappa\) in the accompanying reaction
table is precisely \(\Sigma_n\).

Unlike a ballistic neutron, a gliding cluster performs a 1D random walk and
repeatedly retraces its path. Its absorption time therefore scales as a
squared distance divided by \(D_n\), not as a distance divided by a speed.

\subsection{First-passage derivation of the sink strength}
Let \(p(z,t)\) be the probability density on a sink-free chord bounded by
absorbing interceptions at \(z=0\) and \(z=L\):
\begin{equation}
 \frac{\partial p}{\partial t}=D_n\frac{\partial^2p}{\partial z^2},
 \qquad p(0,t)=p(L,t)=0.                            \label{eq:1d-fp}
\end{equation}
The mean first-passage time \(T(z)\) follows from the backward equation
\begin{equation}
 D_nT''(z)=-1,\qquad T(0)=T(L)=0,
\end{equation}
and hence
\begin{equation}
 T(z)=\frac{z(L-z)}{2D_n}.                          \label{eq:mfpt-position}
\end{equation}
For a uniformly sampled starting point on a chord,
\begin{equation}
 \overline T(L)=\frac1L\int_0^LT(z)\,\dd z
 =\frac{L^2}{12D_n}.                               \label{eq:mfpt-chord}
\end{equation}
These equations establish the fundamental scaling
\(\KK\sim D_n/L^2\sim D_n\Sigma_n^2\).

To resolve the numerical coefficient and the absorbing family, average over
the marked Poisson sequence of line interceptions. An interception is labelled
\(m\) with probability
\begin{equation}
 P(m)=\frac{\Sigma_{n,m}}{\Sigma_n}.
\end{equation}
The self-consistent chemical-rate average over the two bounding
interceptions gives the 1D screening length
\begin{equation}
 \lambda_{s,n}^{-2}=6\Sigma_n^2.                   \label{eq:screening}
\end{equation}
The coefficient differs from the inverse of
\(\overline T(\lambda_g)\) because the physical ensemble contains a
distribution of chord lengths and is weighted by surviving walkers, rather
than identical chords of length \(\lambda_g\). Resolving the terminal
interception by its mark gives
\begin{align}
 k^2_{(1),n,m}
 &=P(m)\lambda_{s,n}^{-2}
 =6\Sigma_{n,m}\Sigma_n \notag\\
 &=6\sigma_{n,m}N_m\sum_j\sigma_{n,j}N_j.
 \label{eq:k1-family}
\end{align}
This is the Borodin--Barashev 1D sink strength
\cite{Borodin1998,Barashev2001}. It follows from line interception, 1D
first passage, and the marked-sink ensemble; it is not a Smoluchowski
\(4\pi RD\) result.

The loss frequency of mobile cluster \(n\) to family \(m\) is
\begin{equation}
 \boxed{\KK^{(1)}_{n\rightarrow m}
 =D_nk^2_{(1),n,m}
 =6D_n\sigma_{n,m}N_m\Sigma_n.}                   \label{eq:K1-Hz}
\end{equation}
In atomic-fraction variables,
\begin{equation}
 \boxed{\KK^{(1)}_{n\rightarrow m}
 =\frac{6D_n\sigma_{n,m}}{\Omega^2}
 c_m\sum_j\sigma_{n,j}c_j
 =24\omega_n\frac{\sigma_{n,m}}{a^4}
 c_m\sum_j\sigma_{n,j}c_j.}                       \label{eq:K1-af}
\end{equation}
The population equation is
\begin{equation}
 \left.\dot c_n\right|_m=-\KK^{(1)}_{n\rightarrow m}c_n.
\end{equation}

\subsection{Origin and consequences of the \(R^4\) law}
For one monodisperse family, with \(\sigma=\pi R^2\),
\begin{equation}
 k^2_{(1)}=6\sigma^2N^2=6\pi^2R^4N^2,\qquad
 \KK^{(1)}=6\pi^2D_nR^4N^2.                       \label{eq:R4}
\end{equation}
The two factors of \(R^2\) have different roles. One cross section selects
the family responsible for terminal capture; the second fixes the spacing
between all line interceptions and hence the distance explored by 1D
diffusion. This produces a loss frequency quadratic in sink density and
proportional to \(R^4\).

For a spherical cavity containing \(m\) vacancies,
\begin{equation}
 R_m=a\left(\frac{3m}{8\pi}\right)^{1/3},
\end{equation}
so
\begin{equation}
 k^2_{(1),m}\propto m^{4/3}N_m^2.
\end{equation}
The result assumes dilute, statistically independent sinks and a fixed glide
direction. Correlated positions, overlap of projected cross sections, finite
reaction probabilities, and orientation-dependent nonspherical sinks require
modification.

\section{1D migration with rotations: the 1DR--0 kernel}
Suppose cluster \(n\) makes 1D flights and changes glide direction after a
characteristic flight length \(\ell_n\). Its trajectory is 1D below
\(\ell_n\) but becomes isotropic at sufficiently large scales. Let
\(k^2_{(1),n,m}\) be given by \cref{eq:k1-family}, and let
\(k^2_{(3),n,m}\) be the conventional 3D sink strength for the same sinks.
The Trinkaus--Heinisch master curve is
\begin{align}
 k^2_{(1R),n,m}&=y_{n,m}k^2_{(1),n,m},             \label{eq:k1r}\\
 y_{n,m}&=\frac12\left(1+\sqrt{1+4x_{n,m}^2}\right),\label{eq:y}\\
 x_{n,m}&=\left[
 \frac{\ell_n^2k^2_{(1),n,m}}{12}
 +\frac{k^4_{(1),n,m}}{k^4_{(3),n,m}}
 \right]^{-1/2}.                                  \label{eq:x}
\end{align}
This \(x\) is the reciprocal of another common convention; equations
\eqref{eq:y} and \eqref{eq:x} must be copied together
\cite{TrinkausHeinisch2002,Barashev2001}. The hertz-valued kernel is
\begin{equation}
 \boxed{\KK^{(1R)}_{n\rightarrow m}
 =D_ny_{n,m}k^2_{(1),n,m}.}                        \label{eq:K1R}
\end{equation}

\subsection{Pure-1D limit}
If \(\ell_n\sqrt{k^2_{(1)}}\gg1\), absorption normally precedes rotation.
Then \(x\to0\), \(y=1+x^2+O(x^4)\), and
\begin{equation}
 \boxed{\KK^{(1R)}\longrightarrow D_nk^2_{(1)}}.
\end{equation}

\subsection{Frequent-rotation, 3D limit}
If the flight length is short compared with the absorption length, the first
term in \cref{eq:x} vanishes and
\begin{equation}
 x\longrightarrow\frac{k^2_{(3)}}{k^2_{(1)}}.
\end{equation}
In the usual dilute-sink regime \(k^2_{(3)}\gg k^2_{(1)}\), hence \(x\gg1\),
\(y=x+1/2+O(x^{-1})\), and
\begin{equation}
 \boxed{\KK^{(1R)}\longrightarrow D_nk^2_{(3)}}.   \label{eq:3d-limit}
\end{equation}
Thus one frequency recovers both limiting sink strengths. The definition of
\(\ell_n\) (mean, projected, or RMS flight length) and the convention used for
the long-time diffusivity must match the cited master curve.

\section{Reactions between two purely 1D-mobile clusters}
Let clusters \(A\) and \(B\) have glide axes \(\bm e_A,\bm e_B\), diffusivities
\(D_A,D_B\), and mutual capture distance \(R=R_A+R_B\). There are three
geometrically distinct cases.

\subsection{Variant I: collinear or coincident glide lines}
When both species sample the same line family, the other mobile population
also sets the interception spacing. Applying the marked-interception result
symmetrically gives the collinear rate density proposed by Adjanor
\cite{Adjanor2022}:
\begin{equation}
 \Rcal_{\parallel}
 =6\pi^2R^4\left(N_A^2N_BD_B+N_B^2N_AD_A\right).  \label{eq:Rparallel}
\end{equation}
Dividing by \(N_A\) gives the loss frequency of an \(A\) cluster:
\begin{equation}
 \boxed{\KK_{A|B}^{\parallel}
 =6\pi^2R^4\left(D_BN_AN_B+D_AN_B^2\right).}       \label{eq:Kparallel}
\end{equation}
Equivalently,
\begin{equation}
 \boxed{\KK_{A|B}^{\parallel}
 =\frac{6\pi^2R^4}{\Omega^2}
 \left(D_Bc_Ac_B+D_Ac_B^2\right).}                \label{eq:Kparallel-af}
\end{equation}
The kinetics is cubic in populations, so this is an effective many-particle
frequency, not a concentration-independent binary constant. Strictly parallel
but distinct lines do not react under perfect 1D motion unless their
separation is smaller than \(R\); lattice-plane occupancy must then be
included.

\subsection{Variant II: nonparallel intersecting glide lines}
Let \(s_A,s_B\) be displacements along two nonparallel lines that intersect.
Near the intersection,
\begin{equation}
 r^2=s_A^2+s_B^2-2s_As_B\cos\theta,
\end{equation}
and the pair probability obeys
\begin{equation}
 \frac{\partial p}{\partial t}
 =D_A\frac{\partial^2p}{\partial s_A^2}
 +D_B\frac{\partial^2p}{\partial s_B^2},
 \qquad p=0\quad(r=R).                             \label{eq:2d-map}
\end{equation}
There are two independent diffusive coordinates, so the encounter maps to a
two-dimensional first-passage problem. For equal diffusivities, the standard
midpoint mapping gives the long-time result
\begin{align}
 \Rcal_{\times}&\simeq K^{\rm vol}_{\times}(t)N_AN_B,\\
 K^{\rm vol}_{\times}(t)&=2\pi DR\,\alpha(t),       \label{eq:Kcross-vol}\\
 \alpha(t)&\simeq\frac{4}{\ln\!\left(4Dt/\pi R^2\right)}.
 \label{eq:alpha}
\end{align}
The slowly varying logarithm is characteristic of 2D diffusion
\cite{GoseleHuntley1975,Adjanor2022}. The hertz-valued loss kernel is
\begin{equation}
 \boxed{\KK_{A|B}^{\times}(t)
 =2\pi DR\,\alpha(t)N_B
 =\frac{2\pi DR}{\Omega}\alpha(t)c_B.}             \label{eq:Kcross-Hz}
\end{equation}
This contribution is approximately quadratic in concentration. For
\(D_A\ne D_B\) or nonorthogonal axes, \cref{eq:2d-map} transforms to diffusion
toward an ellipse. The diffusion-tensor transformation or a validated
anisotropic correlation must be used; replacing \(D\) by \(D_A+D_B\) is not
generally valid \cite{GoseleSeeger1976,Adjanor2022,HuangMarian2019}.

\subsection{Variant III: nonparallel skew, non-coplanar glide lines}
Let \(h\) be the minimum separation of the two infinite glide lines. At their
closest points, contact requires \(h^2+q^2\le R^2\), where \(q\) is distance
in the two-dimensional pair-coordinate plane. Therefore
\begin{equation}
 R_{\eff}(h)=
 \begin{cases}
 \sqrt{R^2-h^2},&h<R,\\
 0,&h\ge R.
 \end{cases}                                      \label{eq:Reff}
\end{equation}
The intersecting-line derivation applies with \(R\to R_{\eff}\):
\begin{equation}
 \boxed{\KK_{A|B}^{\rm skew}(t,h)
 =2\pi D R_{\eff}(h)\,
 \alpha[t;R_{\eff}(h)]N_B.}                        \label{eq:Kskew}
\end{equation}
For a lattice population, average over admissible plane separations and glide
variants. Since \(\alpha\) also depends on radius, the consistent average is
\begin{equation}
 \overline{R_{\eff}\alpha}
 =\sum_hP(h)R_{\eff}(h)\alpha[t;R_{\eff}(h)],
\end{equation}
not necessarily \(\bar R_{\eff}\alpha(t;\bar R_{\eff})\).

\subsection{Orientation-averaged 1D--1D kernel}
If \(v\) unoriented glide-axis variants are equally populated, the fraction of
non-collinear pairs is
\begin{equation}
 f_v=\frac{v-1}{v}.
\end{equation}
For the four unoriented \(\langle111\rangle\) axes in BCC iron, \(f_v=3/4\).
Separating intersecting and skew nonparallel pairs by conditional probability
\(p_\times\) gives
\begin{equation}
 \boxed{\KK_{A|B}^{1D-1D}
 =(1-f_v)\KK_{A|B}^{\parallel}
 +f_v\left[p_\times\KK_{A|B}^{\times}
 +(1-p_\times)\langle\KK_{A|B}^{\rm skew}\rangle\right].}
 \label{eq:total-1d1d}
\end{equation}
This displays what is hidden when only the \(6\pi^2R^4\) term is quoted. In a
dilute system the approximately quadratic non-collinear contribution commonly
exceeds the cubic collinear contribution. Spatial correlations, rotations,
finite cell size, and unequal occupation of glide variants change the weights.

\section{Implementation summary}
For each mobile class \(n\) and fixed family \(m\):
\begin{enumerate}
 \item Compute \(\sigma_{n,m}\) and
 \(\Sigma_n=\sum_j\sigma_{n,j}N_j\).
 \item Form \(k^2_{(1),n,m}=6\sigma_{n,m}N_m\Sigma_n\).
 \item If rotations occur, compute \(x,y\), and
 \(\KK^{(1R)}_{n\to m}=D_nyk^2_{(1),n,m}\).
 \item For two purely 1D-mobile classes, identify the collinear, intersecting,
 or skew geometry and use the corresponding frequency.
 \item Apply orientation and lattice-plane weights before summing frequencies.
 \item For an \(A+A\) reaction, use the conventional factor \(1/2\) in the
 number of events; two \(A\) objects are removed by each event.
\end{enumerate}

\paragraph{Scope.}
The formulas assume diffusion-controlled, perfectly absorbing encounters in a
homogeneous dilute medium. They exclude reaction barriers, elastic drift,
correlated cascade positions, emission, and rotation changes caused by
trapping or stress. Such effects should modify the capture probability,
geometry, or migration model explicitly rather than being hidden in \(R\).

\bibliographystyle{unsrt}
\bibliography{one_dimensional_reaction_kernels}
\end{document}
